\subsection{Uniform recovery and exact wild bands}
Primitive coordinates give a uniform comparison without a restriction on the $p$-part of the absolute ramification index; a compatible cyclic-point construction gives exact mixed thresholds and infinitely many sharp full-class values.

\begin{theorem}[A uniform range at every index]\label{thm:all-index-upper}
Let $p\ge5$ and $n\ge2$. In the standing height-two filtration class, if
\[
                           e_K\le p^{n-1},
\]
then isomorphic actual $p^n$-torsion forces equality of the special inertia characters. This holds over every standing complete field of that index, with no restriction on $v_p(e_K)$. Equivalently,
\[
 \nuMin(p,e,2)\le\left\lceil\log_p e\right\rceil
\]
for every positive multiple $e$ of $p-1$.
\end{theorem}

\begin{theorem}[The mixed threshold at every index]\label{thm:mixed-all-indices}
Let $p\ge5$, $n\ge2$ and let $e$ be a positive multiple of $p-1$. Put $v=v_p(e)$. Over some standing field of exact index $e$, a whole potentially supersingular height-two group and a group filtered by two potentially good height-one groups of total dimension one have isomorphic actual $p^n$-torsion and different special characters if and only if
\[
                              e>p^{2n-3-v}.
\]
The existence field may always have residue field $\overline{\F}_p$. A negative exponent is interpreted in $\Q_{>0}$.
\end{theorem}

\begin{theorem}[Infinite exact wild bands]\label{thm:wild-bands-exact}
For every $p\ge5$, integer $r\ge0$ and integer $2\le k\le(p^2-1)/2$,
\[
 \nuMin\bigl(p,kp^r(p-1),2\bigr)=
 \begin{cases}
 r+2,&p\nmid k\ \text{or }k=p,\\
 r+3,&p\mid k,\ k\ge2p.
 \end{cases}
\]
In the first case, $p^{r+3}$-torsion recovers the special character over every standing field of that exact index, while some such field supports two actual permitted groups with isomorphic $p^{r+2}$-torsion and different special characters. In the second case, these two levels are $p^{r+4}$ and $p^{r+3}$, respectively. In particular,
\[
 \begin{gathered}
 \nuMin\bigl(p,kp(p-1),2\bigr)=3\quad(2\le k\le2p-1),\\
 \nuMin\bigl(p,p^s(p-1),2\bigr)=s+1\quad(s\ge1).
 \end{gathered}
\]
At $k=2p$, this gives
\[
 \nuMin\bigl(p,2p^{r+1}(p-1),2\bigr)=r+3.
\]
\end{theorem}

Section~\ref{sec:one-slope-all-indices} proves the uniform theorem by comparing two Frobenius chains and the larger exact band by matching residual characters and transporting a canonical quotient. Section~\ref{sec:all-wild-indices} gives a separate graph proof of the sufficient range $e\le(p-1)p^{n-2}$, proves the exact mixed criterion and constructs the lower examples for the bands. The mixed criterion determines exactly when an actual group of one reduction type can have the same truncation as a group of the other type. Comparisons between two supersingular groups outside the proved ranges remain a separate question, so the infinite exact bands do not constitute a formula for every later wild index.

